\section{What the Image May and May Not Do}

Used explicitly as poetry, the original image may awaken a true intuition:
Calvary is not sealed away from Christian worship, and the apostles received
more than a prediction. But it should immediately be translated into the
Church's grammar.

\begin{itemize}
  \item Say that the Last Supper \emph{sacramentally anticipated} the Cross;
  do not say that the apostles literally arrived at a future execution.
  \item Say that the Mass \emph{sacramentally makes present} the one sacrifice;
  do not say that the congregation leaves its historical time.
  \item Say that Christ is the same priest and victim; do not imagine the
  risen Lord suffering or dying again.
  \item Say that anamnesis is effective ecclesial memorial; do not reduce it
  either to mental recall or to a hidden physical mechanism.
  \item Say that the Eucharist looks toward the Lord's return and heavenly
  banquet as well as toward the Cross.
  \item Say which statement is doctrine and which is an explanatory model;
  do not present one theologian's mechanism as ``what the Church teaches.''
\end{itemize}

These qualifications do not drain the mystery. They prevent an imaginative
model from becoming a rival explanation. Sacraments are not weak because they
are sacramental. Through created signs and the power of Christ and the Spirit,
they communicate the saving reality in the mode Christ instituted.

